\documentclass[11pt,a4paper]{article}

\usepackage[utf8]{inputenc}
\usepackage[T1]{fontenc}
\usepackage[margin=2.5cm]{geometry}
\usepackage{amsmath,amssymb}
\usepackage{booktabs}
\usepackage{graphicx}
\usepackage{caption}
\usepackage{enumitem}
\usepackage{listings}
\usepackage{xcolor}
\usepackage[hidelinks]{hyperref}
\usepackage{titlesec}

\definecolor{codebg}{RGB}{247,247,250}
\definecolor{codekw}{RGB}{20,90,160}
\definecolor{codecm}{RGB}{110,120,130}

\lstset{
  backgroundcolor=\color{codebg},
  basicstyle=\ttfamily\footnotesize,
  keywordstyle=\color{codekw}\bfseries,
  commentstyle=\color{codecm}\itshape,
  breaklines=true,
  frame=single,
  rulecolor=\color{gray!30},
  language=Python,
  showstringspaces=false,
  xleftmargin=6pt,
  framexleftmargin=6pt
}

\titleformat{\section}{\normalfont\Large\bfseries}{\thesection}{1em}{}
\setlist{itemsep=2pt,topsep=4pt}

\title{\vspace{-1.5cm}\textbf{A Lean Quantum Document Pipeline}\\[0.3cm]
\large Hybrid Quantum Machine Learning and Quantum String Algorithms\\
for Document Extraction, Benchmarked Against Classical Baselines}

\author{Pushkar Kumar\\[0.2cm]
\normalsize QIntern 2026 --- QWorld\\
\normalsize Mentor: Potluri Krishna Priyatham}

\date{August 2026}

\begin{document}
\maketitle

\begin{center}
\small
Repository: \url{https://github.com/Pushkar0997/qi26_25}\\
Live demo: \url{https://pushkar0997.github.io/qi26_25/}
\end{center}

\vspace{0.3cm}

\begin{abstract}
\noindent
Document intelligence pipelines built on large language models carry substantial
computational cost, motivating interest in whether quantum components can
deliver comparable capability at reduced resource requirements. This project
implements an end-to-end document extraction pipeline in which two of nine
stages are quantum: a quanvolutional feature extractor operating on character
crops, and a Grover-based comparator oracle performing pattern search over the
extracted identifier field. Both are benchmarked against dimension-matched
classical equivalents on a controlled synthetic corpus of 60 documents spanning
five degradation tiers.

The central results are negative and are reported as such. Measured across 30
paired train--test splits, the quanvolutional feature extractor is statistically
indistinguishable from a dimension-matched random classical convolution
($p \approx 0.2$--$0.3$, winning approximately half of splits), and both are
outperformed by raw pixel values at lower dimensionality ($p < 10^{-10}$, 30 of
30 splits). Optimising the filter's rotation angles does not close this gap. For
the search stage, a genuine in-circuit comparator oracle is constructed and
verified with zero uncomputation leakage; however, measured two-qubit gate count
scales as $N^{1.39}$ with text length, so the $O(\sqrt{N})$ query advantage of
Grover's algorithm does not survive the cost of loading classical text into the
oracle. Against Tesseract on identical images, the pipeline is three to four
times worse on clean input.

Two methodological findings proved more consequential than any quantum result. A
train/serve distribution mismatch between training and inference crops inflated
document character error rate by a factor of 6.5; correcting it was the single
largest improvement in the project. Separately, isolated-character accuracy
proved an actively misleading proxy for end-to-end accuracy, with a
higher-scoring configuration producing substantially worse documents. The
contribution of this work is therefore a measured characterisation of why the
quantum components did not help at this scale, supported by a fully reproducible
artifact.
\end{abstract}

\vspace{0.3cm}
\tableofcontents
\newpage

%======================================================================
\section{Introduction}

\subsection{Motivation}

Extracting structured information from document images---invoices, forms,
archival records---is a mature application of classical computer vision and
optical character recognition (OCR). Contemporary approaches increasingly route
this task through large language models, which achieve strong accuracy at
considerable computational expense. The QIntern 2026 project brief posed the
question of whether quantum machine learning and quantum string algorithms could
deliver a ``lean'' alternative: comparable extraction capability at reduced
resource cost.

This question decomposes naturally into two independent sub-problems, which the
brief assigned to separate tracks:

\begin{itemize}
\item \textbf{Track A (vision).} Can a quantum feature extractor---specifically a
      quanvolutional layer, the quantum analogue of a convolution---provide
      useful features for character recognition?
\item \textbf{Track B (search).} Can Grover's algorithm accelerate pattern search
      over extracted text, exploiting its $O(\sqrt{N})$ query complexity?
\end{itemize}

Both are legitimate questions with published affirmative-leaning literature, and
both are answered negatively here, for reasons that are measured rather than
asserted.

\subsection{Scope and Honest Framing}

This work makes no claim of quantum advantage. The explicit goal, following the
project brief, was to construct a working pipeline and characterise its resource
requirements against classical baselines. Where the quantum components
underperform, that is reported directly, together with the mechanism
responsible.

Three constraints bound what can be concluded. First, all quantum execution is
\emph{simulated}; no physical quantum hardware was available. Second, the corpus
is synthetic, which is necessary for exact ground truth but does not represent
real scanned documents. Third, the scale is small---4 qubits, $8\times8$ crops,
a 36-class alphabet---so conclusions apply to this regime and should not be
extrapolated.

\subsection{Contributions}

\begin{enumerate}
\item A complete nine-stage pipeline in which quantum feature extraction and
      quantum search are integrated with classical segmentation and
      classification (Section~\ref{sec:pipeline}).
\item A comparator-based Grover oracle that performs string comparison
      \emph{inside} the circuit, replacing a common construction that computes
      match positions classically before circuit construction
      (Section~\ref{sec:oracle}).
\item Statistical evaluation of the quantum-versus-classical feature comparison
      across 30 paired splits, distinguishing a null result from a real effect
      (Section~\ref{sec:features}).
\item Measured resource scaling for the Grover oracle, demonstrating that data
      loading dominates the quadratic query advantage (Section~\ref{sec:scaling}).
\item An external baseline against Tesseract, and identification of the
      classical segmentation front end---not either quantum stage---as the
      dominant error source (Section~\ref{sec:endtoend}).
\item A fully reproducible artifact: committed dataset, committed model,
      continuous integration that fails if the report and the measured results
      disagree (Section~\ref{sec:repro}).
\end{enumerate}

%======================================================================
\section{Background}

\subsection{Quantum Image Encoding}

Before any quantum processing can occur, pixel values must become quantum states.
Two standard encodings were evaluated.

\textbf{FRQI} (Flexible Representation of Quantum Images) stores intensity as a
rotation angle on a single shared colour qubit:
\begin{equation}
|I\rangle = \frac{1}{2^{n}}\sum_{i=0}^{2^{2n}-1}
\big(\cos\theta_i\,|0\rangle + \sin\theta_i\,|1\rangle\big)\otimes|i\rangle
\end{equation}
One qubit carries all colour information regardless of image size, but intensity
is recovered only statistically: $\theta$ must be estimated from measurement
frequencies, so reconstruction is approximate and shot-limited.

\textbf{NEQR} (Novel Enhanced Quantum Representation) stores intensity as a bit
string on a colour register:
\begin{equation}
|I\rangle = \frac{1}{2^{n}}\sum_{i}|f(i)\rangle\otimes|i\rangle
\end{equation}
An 8-bit greyscale image requires 8 colour qubits, but reconstruction is exact.

\subsection{Quanvolutional Layers}

Introduced by Henderson et al., a quanvolutional layer is the quantum analogue of
a convolution: a small window slides over the image and, at each position, the
patch is encoded into a quantum state, transformed by a circuit, and measured.
The canonical formulation uses a \emph{fixed random} circuit---the filter is not
trained---which is the configuration evaluated here, with a trained variational
variant examined separately in Section~\ref{sec:variational}.

\subsection{Grover Search and String Matching}

Grover's algorithm locates a marked item among $N$ in $O(\sqrt{N})$ oracle
queries. Applying it to string matching promises quadratically faster pattern
location. The critical requirement, examined in detail in
Section~\ref{sec:oracle}, is that the oracle must \emph{perform} the comparison
rather than being supplied its result.

%======================================================================
\section{Pipeline Architecture}
\label{sec:pipeline}

The pipeline comprises nine stages, of which stages 4--5 and 9 are quantum:

\begin{center}
\small
\begin{tabular}{@{}rll@{}}
\toprule
\# & Stage & Description \\
\midrule
1 & preprocess & median filter, Otsu threshold, speckle removal \\
2 & segment & projection profiles $\rightarrow$ character boxes \\
3 & crop & pad to square, resize to $8\times8$, normalise \\
4 & \textbf{encode} & \textbf{$R_Y$ angles $\rightarrow$ 4 qubits per patch (quantum)} \\
5 & \textbf{quanvolutional} & \textbf{entangling circuit $\rightarrow$ 160 features (quantum)} \\
6 & classify & multinomial logistic regression $\rightarrow$ characters \\
7 & assemble & rebuild lines, insert spaces \\
8 & locate & identify the document identifier field \\
9 & \textbf{Grover search} & \textbf{comparator oracle $\rightarrow$ pattern position (quantum)} \\
\bottomrule
\end{tabular}
\end{center}

\subsection{Dataset}

Sixty documents were generated across five degradation tiers---clean digital,
clean scan, noisy scan, clear handwriting, degraded handwriting---yielding
\textbf{6{,}023 labelled character crops} over 36 classes. Generation is
deterministic under seed 26.

Synthetic generation was chosen for two reasons. Character error rate requires
exact ground truth, and hand-labelling scraped documents was not feasible within
the project timeframe. More importantly, every tier renders the \emph{same
distribution of source text}, so any accuracy difference between tiers is
attributable to image degradation alone rather than to differing content
difficulty.

The degradation is a controlled variable: blur, sensor noise, contrast
reduction, low-DPI resampling, rotation and blotch artifacts, with per-tier
parameters recorded in the generator. That the tiers form a genuine difficulty
gradient is confirmed by median per-crop dynamic range, which falls
monotonically from 197 (clean digital) to 62 (degraded handwriting).

\paragraph{Limitations of the corpus.} The handwriting tiers use an italic serif
face with per-glyph position, rotation and scale jitter as a proxy for
handwriting variability. This is not handwriting, and the IAM Handwriting
Database is identified as the correct extension. This limitation is stated in
the generator source, the repository README and here.

%======================================================================
\section{Track A: Quantum Feature Extraction}

\subsection{Encoding Comparison: FRQI versus NEQR}
\label{sec:encoding}

Both encodings were implemented, verified by reconstruction, and measured across
three patch sizes.

\begin{table}[h]
\centering
\caption{Encoding resource comparison. Gate counts are measured after
transpilation to a $\{u, \mathrm{CX}\}$ basis.}
\begin{tabular}{@{}lrrrrr@{}}
\toprule
Patch & FRQI qubits & NEQR qubits & FRQI CX & NEQR CX & Qubit ratio \\
\midrule
$2\times2$ & 3 & 10 & 32 & 114 & $3.33\times$ \\
$4\times4$ & 5 & 12 & 384 & 1{,}558 & $2.40\times$ \\
$8\times8$ & 7 & 14 & 3{,}584 & 10{,}998 & $2.00\times$ \\
\bottomrule
\end{tabular}
\end{table}

Verification confirmed the expected behaviour: NEQR reconstructs pixel values
exactly, while FRQI reconstructs to within 1--2 intensity levels out of 255 at
8{,}192 shots.

\paragraph{A refuted prediction.} An early technical summary written during this
project predicted that the qubit gap between the two encodings would
\emph{widen} with patch size. Measurement refutes this: the ratio
\emph{narrows}, from $3.33\times$ to $2.00\times$. The original reasoning
conflated two quantities. FRQI's colour register is fixed at 1 qubit and NEQR's
at 8, so the colour gap is a \emph{constant} 7 qubits at every patch size, while
both encodings share an identically growing position register of
$\lceil\log_2(\text{pixels})\rceil$. A constant gap divided by a growing shared
base necessarily shrinks as a ratio.

Absolute CX counts grow steeply for both, with NEQR reaching approximately
11{,}000 at $8\times8$. It is specifically the \emph{relative} qubit penalty that
decays. The pipeline uses neither encoding directly---a patch-based
quanvolutional filter does not require position addressing---but these
measurements establish the resource baseline for whole-image quantum
representation, and the $8\times8$ NEQR figure is the quantitative reason the
pipeline is patch-based.

\subsection{The Quanvolutional Filter}

The filter is a fixed 4-qubit circuit: two layers of $R_Y$ rotations with a ring
of CX gates between them. Rotation angles are drawn once from a seeded generator
and held constant.

Each $8\times8$ character crop is divided into sixteen $2\times2$ patches. Each
pixel value $v$ becomes a rotation angle $\theta = \pi v/255$ applied to its own
qubit, so a dark pixel remains near $|0\rangle$ and a bright one rotates toward
$|1\rangle$. Because each qubit is rotated independently, the input is a
\emph{product state}; all entanglement in the output originates in the filter's
CX gates.

\subsubsection{Efficient Simulation}

A naive implementation submits one simulator job per patch, requiring
approximately 96{,}000 jobs for a single pass over the dataset---several hours of
wall-clock time, which makes iteration impractical.

Because the filter is fixed, its unitary $U$ is identical for every patch. It is
extracted once and applied as a batched matrix multiplication:
\begin{equation}
|\psi_{\text{out}}\rangle = U\,|\psi_{\text{in}}\rangle
\end{equation}
This is not an approximation but the same linear algebra the simulator performs,
without per-job overhead and without shot noise. Agreement with the shot-based
implementation was verified at a maximum deviation of \textbf{0.0028} across
random patches at 200{,}000 shots---within sampling error. Processing time fell
to 0.15 seconds for 2{,}000 characters.

\subsubsection{Readout Design}

The obvious readout is per-qubit $P(\text{measure }1)$: four values. However, a
4-qubit circuit produces a distribution over 16 outcomes, and retaining only four
marginals discards every correlation the entangling layer created---precisely the
component a classical product-state model could not produce.

Adding pairwise correlations $\langle Z_i Z_j\rangle$ recovers six additional
values per patch, giving 160 features per character rather than 64.

\begin{table}[h]
\centering
\caption{Effect of readout design, all tiers pooled.}
\begin{tabular}{@{}lrr@{}}
\toprule
Readout & Feature dimension & Accuracy \\
\midrule
Marginals only & 64 & 86.7\% \\
Marginals $+$ $\langle Z_i Z_j\rangle$ & 160 & \textbf{92.2\%} \\
\bottomrule
\end{tabular}
\end{table}

The 5.5-point gain from readout design alone exceeds any difference measured
between the quantum and classical filters, which is itself a result: for
quanvolutional approaches, how much of the output distribution is retained
matters more than which circuit produced it.

%======================================================================
\section{Track B: Quantum String Matching}
\label{sec:oracle}

\subsection{A Construction That Does Not Search}

A common formulation of Grover-based string matching, and the one present in this
project's own starter notebook, proceeds as follows:

\begin{lstlisting}
match_positions = [i for i in range(len(text) - len(pattern) + 1)
                   if text[i:i+len(pattern)] == pattern]
oracle = build_marked_state_oracle(n_qubits, match_positions)
\end{lstlisting}

The comparison executes in Python, before any qubit is allocated. The circuit is
then constructed to flip the phase of positions already known. This demonstrates
Grover \emph{amplification} correctly---the measured distribution peaks at the
correct position---but it is not string matching, and no complexity claim can
rest on it, because the $O(N)$ classical scan has already occurred.

\subsection{An In-Circuit Comparator Oracle}

The replacement receives only the text and the pattern. Two registers are used:
a position register \texttt{pos} of $\lceil\log_2(N-M+1)\rceil$ qubits holding
candidate start positions in superposition, and a window register \texttt{win} of
$M \times b$ qubits, where $b$ is the alphabet width in bits.

Four stages execute:

\begin{enumerate}
\item \textbf{Window loading.} For each concrete position $i$, multi-controlled
      writes load the bits of $\text{text}[i:i+M]$ into \texttt{win},
      conditioned on $\texttt{pos} = i$. Afterwards
      $|i\rangle|0\rangle \rightarrow |i\rangle|\text{text}[i:i+M]\rangle$.
\item \textbf{Pattern XOR.} An $X$ gate is applied to $\texttt{win}[j]$ wherever
      pattern bit $j$ is 1. The register is now all-zeros exactly when the
      window equals the pattern.
\item \textbf{Phase marking.} A multi-controlled phase gate, sandwiched between
      $X$ layers, flips the phase iff \texttt{win} is all-zeros.
\item \textbf{Uncomputation.} Stages 2 and 1 are reversed, disentangling
      \texttt{win} and returning it to $|0\rangle$.
\end{enumerate}

Stage 4 is not housekeeping. Without it, the window register remains entangled
with the position register and functions as a which-path record: the amplitudes
cannot interfere and the diffusion operator amplifies nothing.

\paragraph{An encoding detail.} Alphabet codes begin at 1, never 0. When
$N-M+1$ is not a power of two, the position register addresses more states than
there are valid positions; those unreachable states never trigger a controlled
load, so their window remains all-zeros---exactly the pattern stage 3 marks as a
match. Reserving code 0 eliminates this failure mode.

\subsection{Staged Verification}

Correctness was established in stages rather than by observing a plausible final
output.

\begin{table}[h]
\centering
\caption{Oracle verification. Leakage is total probability outside the
$\texttt{win}=0$ subspace after uncomputation.}
\begin{tabular}{@{}llc@{}}
\toprule
Stage & Check & Result \\
\midrule
1 & Window register holds $\text{text}[i:i+M]$ for a probe position & pass \\
2 & Exactly the true match positions have negative amplitude & pass \\
3 & Uncomputation leakage & $\mathbf{0.00}$ \\
4 & Full search returns a true match & pass, 95--96\% confidence \\
\bottomrule
\end{tabular}
\end{table}

On a 16-character hexadecimal field with a 2-character pattern: 14 qubits, 3
Grover iterations, correct position recovered in 95.2\% of shots.

\subsection{Resource Scaling}
\label{sec:scaling}

Grover's algorithm provides $O(\sqrt{N})$ \emph{queries}. However, the window
loading stage must write the text into the circuit, and each candidate position
requires its own controlled load.

\begin{table}[h]
\centering
\caption{Measured oracle resources against text length, hexadecimal alphabet,
2-character pattern.}
\begin{tabular}{@{}rrrr@{}}
\toprule
$N$ (characters) & Qubits & Depth & CX gates \\
\midrule
8 & 13 & 1{,}857 & 1{,}082 \\
12 & 14 & 4{,}583 & 2{,}492 \\
16 & 14 & 6{,}383 & 3{,}428 \\
24 & 15 & 11{,}445 & 6{,}122 \\
32 & 15 & 16{,}159 & 8{,}570 \\
48 & 16 & 29{,}664 & 15{,}696 \\
64 & 16 & 39{,}558 & 20{,}904 \\
\bottomrule
\end{tabular}
\end{table}

Fitting a power law in log--log space gives a CX scaling exponent of
$\mathbf{1.39}$: \emph{superlinear} in text length. The measured curve lies above
a linear reference, nowhere near $\sqrt{N}$.

\paragraph{Interpretation.} For searching stored, unstructured classical text
that must be loaded into the circuit, the quadratic query advantage does not
translate into an end-to-end speedup; the loading cost dominates. This is a
well-known caveat in the literature, and the contribution here is reproducing it
with direct measurement rather than quoting $\sqrt{N}$ unqualified.

Grover's advantage remains real where the text is already available as a quantum
oracle or in QRAM; where the search space is \emph{generated} by a function
rather than stored, as in satisfiability or cryptographic preimage search; or
where a loaded state is reused across many queries, amortising the cost. None of
these describes scanning a document just received.

A second practical constraint: a full 38-symbol alphabet requires 6 qubits per
character, placing even a short window beyond feasible simulation. The pipeline
therefore searches the hexadecimal-representable identifier field, which is the
realistic target for structured-field lookup in any case.

%======================================================================
\section{Benchmarking}
\label{sec:features}

\subsection{Experimental Design}

All feature extractors were compared \emph{at matched dimensionality}, on
identical crops, with an identical classifier and split. Without matching
dimensionality, a win would be indistinguishable from granting one method a wider
representation.

\begin{table}[h]
\centering
\caption{Character accuracy by extractor and tier. Dimensions in parentheses.}
\small
\begin{tabular}{@{}lrrrr@{}}
\toprule
Tier & quanv (64) & quanv$+ZZ$ (160) & classical (160) & raw pixels (64) \\
\midrule
clean\_digital & 99.3\% & 100\% & 100\% & 100\% \\
clean\_scan & 99.3\% & 100\% & 100\% & 100\% \\
noisy\_scan & 93.0\% & 97.7\% & 97.7\% & 98.3\% \\
clear\_handwriting & 87.7\% & 96.7\% & 96.3\% & 98.7\% \\
degraded\_handwriting & 62.5\% & 75.7\% & 74.4\% & 79.4\% \\
\midrule
\textbf{All tiers} & \textbf{86.7\%} & \textbf{92.2\%} & \textbf{91.9\%} & \textbf{93.6\%} \\
\bottomrule
\end{tabular}
\end{table}

\subsection{Statistical Significance}

A single split yields point estimates but cannot distinguish a real difference
from split-to-split noise. The comparison was repeated over 30 splits, with every
extractor scored on the same splits so that comparisons are paired.

\begin{table}[h]
\centering
\caption{Paired comparisons across 30 splits. Figures are given as ranges because
scikit-learn's solver is not bit-reproducible across BLAS builds; conclusions are
stable across runs, individual decimal places are not.}
\small
\begin{tabular}{@{}lrrl@{}}
\toprule
Comparison & Mean difference & $p$ & Verdict \\
\midrule
quanv$+ZZ$ $-$ classical & $-0.1$ pt & $\approx 0.2$--$0.3$ & not distinguishable from noise \\
raw pixels $-$ quanv$+ZZ$ & $+1.2$ pt & $< 10^{-10}$ & significant \\
raw pixels $-$ classical & $+1.1$ pt & $< 10^{-12}$ & significant \\
quanv$+ZZ$ $-$ quanv & $+5.7$ pt & $< 10^{-25}$ & significant \\
\bottomrule
\end{tabular}
\end{table}

This separates two claims that a single-split table presents identically. The
quantum-versus-classical gap is noise: the quantum layer wins approximately half
of the splits. The raw-pixel lead is not: it holds in all 30.

\subsection{Interpretation}

Three findings, in order of importance.

\textbf{At matched dimensionality, the quantum layer performs at parity with its
classical analogue.} It is not \emph{specifically} worse; a random classical
convolution with the same patch size, stride and output width scores the same
within noise. Whatever limits performance is not quantumness.

\textbf{Both lose to raw pixels, at lower dimensionality.} This reframes the
comparison: the relevant axis is not quantum versus classical but \emph{untrained
random feature map} versus \emph{no feature map}. At $8\times8$ the input is
already close to information-minimal, so any fixed random projection discards
more than it contributes. That the classical control fails identically is what
identifies the cause.

\textbf{Readout design dominated filter design}, as noted in
Section~\ref{sec:encoding}.

\subsection{Variational Filter}
\label{sec:variational}

The project brief specified a \emph{variational} circuit. The filter's eight
$R_Y$ angles were optimised by gradient-free Powell search, starting from the
untrained baseline's exact angle vector. The classical control was trained
identically---training only the quantum side would invert the fairness the
untrained comparison was designed to preserve.

\begin{table}[h]
\centering
\caption{Variational training, evaluated on a sealed test partition of 1{,}807
crops.}
\begin{tabular}{@{}lrrr@{}}
\toprule
Configuration & Parameters & Accuracy & $\Delta$ \\
\midrule
Quantum, untrained & --- & 91.7\% & --- \\
Quantum, trained & 8 & 90.6\% & $-1.1$ pt \\
Classical, untrained & --- & 92.0\% & --- \\
Classical, trained & 40 & 92.1\% & $+0.1$ pt \\
Raw pixels & --- & \textbf{93.6\%} & --- \\
\bottomrule
\end{tabular}
\end{table}

Training does not close the gap. The Section~\ref{sec:features} conclusion
therefore strengthens: the deficit is not caused by the filters being untrained,
because training them does not help.

\paragraph{A methodological correction.} An earlier version of this experiment
reported a $+1.4$-point gain for the trained quantum filter, bringing it level
with raw pixels. That result was an artifact. The angle search had been run over
the full dataset, and the final evaluation then used a fresh random split of that
same dataset, so the angles had been selected using rows that later appeared in
the test set. Re-running with a sealed 30\% partition the search never sees
reduced the gain from $+1.4$ points to approximately zero. The inflated figure is
recorded here rather than discarded, because the difference between the two is
the difference between a false headline claim and the actual finding.

One property does survive: the quantum filter produces its 160-dimensional
feature space from 8 parameters against the classical control's 40. At the
untrained baseline the two perform equivalently, so that $5\times$ parameter
economy is real and costs nothing there. It does not survive this training
protocol, and should not be reported without that qualification.

\subsection{Shot-Noise Sensitivity}

Exact statevector results are achievable only in simulation; hardware samples.

\begin{table}[h]
\centering
\caption{Accuracy against measurement shot budget per patch.}
\begin{tabular}{@{}lr@{}}
\toprule
Shots & Accuracy \\
\midrule
64 & 70.7\% \\
256 & 81.7\% \\
1{,}024 & 85.2\% \\
4{,}096 & 86.7\% \\
exact & 86.7\% \\
\bottomrule
\end{tabular}
\end{table}

Approximately 1{,}024 shots per patch are required to approach the noiseless
result. At 16 patches per character this is roughly 16{,}000 circuit executions
per character---a substantial hardware cost that belongs in any honest resource
accounting and is frequently omitted from published quanvolutional results.

%======================================================================
\section{End-to-End Results}
\label{sec:endtoend}

Two measures are used. \textbf{Character error rate (CER)} is the edit distance
between extracted and reference text, normalised by reference length.
\textbf{Document ID recovered exactly} is the fraction of documents whose
identifier field is completely correct---the pipeline's actual task, for which a
single wrong character renders the result useless.

\begin{table}[h]
\centering
\caption{End-to-end performance, all 12 documents per tier.}
\begin{tabular}{@{}lrrr@{}}
\toprule
Tier & CER & Segmentation ratio & ID recovered exactly \\
\midrule
clean\_digital & 6.6\% & 0.95 & \textbf{100\%} \\
clean\_scan & 8.5\% & 0.93 & 58\% \\
noisy\_scan & 79.5\% & 0.34 & 0\% \\
clear\_handwriting & 51.8\% & 0.70 & 0\% \\
degraded\_handwriting & 85.7\% & 0.99 & 0\% \\
\bottomrule
\end{tabular}
\end{table}

The two metrics diverge sharply. \texttt{clean\_scan} reads at 8.5\% CER---
visually almost clean---yet only 58\% of its documents yield a usable identifier.
Averaged character accuracy flatters a system that fails the actual task in
almost half of cases.

\subsection{Train/Serve Skew: The Largest Single Error Source}

The largest improvement in this project came from a data-handling defect rather
than any modelling change.

\paragraph{Detection.} Switching the feature stage from the 64-dimensional
marginal readout to the 160-dimensional $+ZZ$ readout \emph{raised} isolated
character accuracy yet made end-to-end CER \emph{worse}. A strictly better
feature set producing strictly worse output is impossible if both are measured on
the same input distribution---so they were not.

\paragraph{Cause.} Training crops were cut from exact glyph bounds recorded
during dataset generation; inference crops come from projection segmentation.
These are measurably different distributions: on \texttt{clean\_scan},
ground-truth boxes have median size $16\times19$ px against segmented boxes at
$12\times15$ px. The higher-capacity features fitted the training distribution
more tightly and transferred less well.

\paragraph{Correction.} Training crops are now built by running the real
segmentation path over training pages and labelling each detected box from ground
truth by overlap, with ambiguous detections dropped rather than guessed.

\begin{table}[h]
\centering
\caption{Effect of correcting train/serve skew.}
\begin{tabular}{@{}lrrr@{}}
\toprule
Tier & CER before & CER after & Reduction \\
\midrule
clean\_digital & 42.9\% & \textbf{6.6\%} & $6.5\times$ \\
clean\_scan & 29.7\% & \textbf{8.5\%} & $3.5\times$ \\
clear\_handwriting & 74.5\% & \textbf{51.8\%} & $1.4\times$ \\
\bottomrule
\end{tabular}
\end{table}

\paragraph{A caution on the accuracy figure.} Training on segmented crops
\emph{lowers} reported held-out character accuracy from 92.6\% to 67.7\%. These
numbers are not comparable---they are measured on different crop
distributions---and the lower one is measured on the harder, realistic one. The
comparison to draw is that a model scoring 92.6\% on clean bounds produced 42.9\%
document CER, while a model scoring 67.7\% on realistic bounds produced 6.6\%.

\textbf{Isolated-character accuracy was an actively misleading proxy for
end-to-end accuracy, and optimising it drove the pipeline in the wrong
direction.} Any comparison of feature extractors on pre-cut crops---including
Section~\ref{sec:features} of this report---measures something narrower than
pipeline quality.

\subsection{External Baseline}

Sections~\ref{sec:features} and~\ref{sec:endtoend} compare feature maps that all
share this project's segmentation front end and classifier. Tesseract provides an
external reference, run on the same images over the same twelve documents per
tier, scored with the same error function.

\begin{table}[h]
\centering
\caption{Comparison against Tesseract.}
\begin{tabular}{@{}lrrrr@{}}
\toprule
Tier & This pipeline & Tesseract & Pipeline ID & Tesseract ID \\
\midrule
clean\_digital & 6.6\% & \textbf{2.0\%} & \textbf{100\%} & 92\% \\
clean\_scan & 8.5\% & \textbf{1.9\%} & 58\% & \textbf{83\%} \\
noisy\_scan & \textbf{79.5\%} & 100.0\% & 0\% & 0\% \\
clear\_handwriting & 51.8\% & \textbf{20.2\%} & 0\% & \textbf{17\%} \\
degraded\_handwriting & \textbf{85.7\%} & 100.0\% & 0\% & 0\% \\
\bottomrule
\end{tabular}
\end{table}

Tesseract is three to four times better on clean and moderately degraded input.
This is the expected result and is reported plainly: four decades of engineering
on a task-specific problem outperforms a 4-qubit filter over $8\times8$ crops.

Two details complicate rather than rescue the picture. On \texttt{clean\_digital}
this pipeline recovers the identifier in every case against Tesseract's 92\%,
because a constrained 36-class charset cannot emit the out-of-charset characters
Tesseract occasionally produces. On the two most degraded tiers Tesseract returns
essentially nothing (100\% error) where this pipeline still recovers roughly one
character in five---an artifact of Tesseract's page-analysis stage rejecting
input it judges unreadable, rather than evidence of a superior method.

The comparison is not controlled: Tesseract brings its own segmentation, language
model and training corpus. It is included because the first question a sceptical
reader asks is how the pipeline compares to an ordinary classical tool, and the
answer belongs in the report regardless of whether it flatters the work.

\subsection{Where the Remaining Error Lies}

On clean tiers, segmentation recovers 93--95\% of characters and isolated
character accuracy approaches 99\%, yet CER is 6.6--8.5\%. The residual is
therefore not classification but missed or merged glyphs and word-spacing
reconstruction. \textbf{The classical front end, not either quantum stage,
is the accuracy bottleneck.}

%======================================================================
\section{Noise and Hardware Feasibility}

All results above assume noiseless simulation. Three degradation mechanisms were
modelled to estimate hardware behaviour: shot noise (Section~\ref{sec:features}),
gate error, and readout error.

Gate noise was modelled with a depolarising channel. Because the filter circuit
is small, the noisy channel was constructed once as a $256\times256$
superoperator and applied by batched matrix multiplication, verified to reproduce
the noiseless path exactly at $p=0$.

\paragraph{A measurement subtlety.} Accuracy barely moves under gate noise when
the classifier is refitted at each noise level. This is misleading: depolarising
noise contracts features uniformly toward the maximally mixed state, and
refitting simply learns larger weights. A flat curve could therefore mean ``the
noise did not matter'' \emph{or} ``the retraining absorbed it.'' Reporting a
train-clean/deploy-noisy transfer measurement alongside removes the ambiguity;
the two agree below $3\times10^{-2}$ error and diverge above it, which is what
establishes the robustness as real rather than an artifact of retraining.

\begin{table}[h]
\centering
\caption{Circuit size against error tolerance, at a two-qubit gate error rate of
$10^{-3}$.}
\begin{tabular}{@{}lrr@{}}
\toprule
Circuit & CX gates & $P(\text{error-free})$ \\
\midrule
Quanvolutional filter, per patch & 8 & 99.2\% \\
Grover oracle, 16-character field & 3{,}272 & 3.8\% \\
Grover oracle, 3 iterations & 9{,}816 & $5\times10^{-5}$ \\
\bottomrule
\end{tabular}
\end{table}

The conclusion is stark and is an independent argument from the data-loading
result of Section~\ref{sec:scaling}: \textbf{the feature extractor is plausible
on near-term hardware; the search stage is not.} The difference is three orders
of magnitude in circuit size. The search stage requires error correction, not
improved calibration.

%======================================================================
\section{Reproducibility}
\label{sec:repro}

Every figure in this report regenerates from a clean repository clone. Three
measures support this.

\textbf{The dataset and trained model are committed.} Both were initially
gitignored as derived artifacts. This was incorrect: generation renders text with
whatever fonts the host machine provides, and training is not bit-reproducible
across BLAS builds. During the project, a regeneration on a machine whose Pillow
was built with libraqm produced identical text with every glyph shifted by 0--2
pixels---enough to change clean-document CER from 6.6\% to 63.7\% in an earlier
configuration. The generator now pins the text layout engine explicitly and the
dataset manifest records the Pillow, FreeType and libraqm state.

\textbf{The browser demonstration is verified against the Python
implementation.} An accompanying interactive demonstration runs the full pipeline
client-side. It is not a re-implementation: it loads the exact exported weights,
and a parity harness asserts stage-by-stage agreement across three quality tiers,
with features matching to $5.6\times10^{-16}$ and decoded text identical.
Achieving this required matching Pillow's antialiased triangle resampling filter,
its integer truncation semantics, its 22-bit fixed-point accumulation, its
integer luma conversion, and SciPy's 4-connectivity default---each of which
silently altered decoded characters when mismatched.

\textbf{Continuous integration gates the claims.} On every push, a workflow
re-runs the benchmark against tolerance bands, executes the oracle self-test,
verifies browser parity, checks that the significance conclusions still hold, and
fails the build if any figure in the results file does not appear in this report.
That final gate has caught a real contamination: a results file measured on a
drifted dataset.

%======================================================================
\section{Conclusions}

\begin{enumerate}
\item A complete hybrid pipeline was constructed and runs end to end: quantum
      feature extraction feeding a classical OCR backend, feeding a genuine
      Grover comparator search over the recovered identifier field.

\item At matched dimensionality the quantum feature layer is statistically
      indistinguishable from a classical convolution ($p \approx 0.2$--$0.3$,
      approximately half of 30 splits), and both are outperformed by raw pixels
      ($p<10^{-10}$, 30 of 30). On this task, at this scale, quantum feature
      extraction provides no accuracy advantage.

\item Grover string matching functions correctly, but measured gate cost grows
      superlinearly (exponent $1.39$) with text length owing to data loading, so
      the $\sqrt{N}$ query advantage does not translate into an end-to-end
      speedup for stored classical text.

\item The dominant error source in the full pipeline is the classical
      segmentation front end, not either quantum stage. Clean-document CER is
      6.6--8.5\%, with 100\% identifier recovery on clean digital documents.

\item The largest single improvement came from correcting a train/serve skew,
      not from any modelling change---and was found only because a better
      feature set produced worse end-to-end output.

\item An early prediction made within this project, that the FRQI/NEQR qubit gap
      would widen with patch size, is refuted by measurement; it narrows from
      $3.33\times$ to $2.00\times$.

\item Against Tesseract on identical images, this pipeline is three to four times
      worse on clean input. It is not competitive with mature classical OCR.
\end{enumerate}

This work does not demonstrate quantum advantage for document intelligence, and
identifies specific measured reasons why: untrained random feature maps
underperform the identity at small input sizes, and oracle data loading erases
Grover's asymptotic edge. Both are useful negative results, and both are more
informative than a tuned demonstration would have been.

\subsection{Future Work}

\begin{enumerate}
\item \textbf{A trained filter with a smooth objective.} The variational
      experiment optimised a step-function accuracy metric with a gradient-free
      method. Parameter-shift gradients on a smooth surrogate loss, with
      cross-validated folds, is the correct approach.
\item \textbf{Replacing the segmentation front end.} It is the measured
      bottleneck; a learned detector would likely improve end-to-end results more
      than any change to the quantum stages.
\item \textbf{Real handwriting.} The handwriting tiers are a font-plus-jitter
      proxy. The IAM Handwriting Database is the appropriate extension.
\item \textbf{Stem-level or block-level encodings for search.} The
      one-qubit-per-position encoding is what makes the oracle expensive;
      coarser encodings would reduce both qubit count and loading cost.
\item \textbf{Execution on physical hardware} for the feature extractor, which
      the noise analysis suggests is feasible.
\end{enumerate}

%======================================================================
\section*{Scope and Contributions}
\addcontentsline{toc}{section}{Scope and Contributions}

This project was scoped for three participants: Klasik Taidi on Track A (vision),
Hoang Dinh Duy Anh on Track B (search), and Pushkar Kumar on orchestration and
integration, mentored by Potluri Krishna Priyatham.

Mentor contact ended in early July. The two track assignments were not taken up,
and neither teammate appears in the repository commit history. Both tracks were
consolidated into a single-person scope, which is why the majority of commits
fall within the final days of the schedule.

%======================================================================
\begin{thebibliography}{9}

\bibitem{henderson}
M. Henderson, S. Shakya, S. Pradhan, T. Cook,
\emph{Quanvolutional Neural Networks: Powering Image Recognition with Quantum
Circuits}, Quantum Machine Intelligence, vol. 2, no. 1, 2020.

\bibitem{grover}
L. K. Grover,
\emph{A fast quantum mechanical algorithm for database search},
Proceedings of the 28th Annual ACM Symposium on Theory of Computing, 1996,
pp. 212--219.

\bibitem{frqi}
P. Q. Le, F. Dong, K. Hirota,
\emph{A flexible representation of quantum images for polynomial preparation,
image compression and processing operations},
Quantum Information Processing, vol. 10, no. 1, pp. 63--84, 2011.

\bibitem{neqr}
Y. Zhang, K. Lu, Y. Gao, M. Wang,
\emph{NEQR: a novel enhanced quantum representation of digital images},
Quantum Information Processing, vol. 12, no. 8, pp. 2833--2860, 2013.

\bibitem{bbht}
M. Boyer, G. Brassard, P. H\o yer, A. Tapp,
\emph{Tight bounds on quantum searching},
Fortschritte der Physik, vol. 46, pp. 493--505, 1998.

\bibitem{qiskit}
Qiskit Contributors,
\emph{Qiskit: An Open-source Framework for Quantum Computing}, 2023.

\bibitem{tesseract}
R. Smith,
\emph{An Overview of the Tesseract OCR Engine},
Proceedings of the Ninth International Conference on Document Analysis and
Recognition (ICDAR), 2007, pp. 629--633.

\bibitem{otsu}
N. Otsu,
\emph{A Threshold Selection Method from Gray-Level Histograms},
IEEE Transactions on Systems, Man, and Cybernetics, vol. 9, no. 1,
pp. 62--66, 1979.

\end{thebibliography}

\end{document}
